% !TeX program = pdflatex
% Zálohování a obnova dat — Tappka
% Sazba: latexmk -pdf zalohovani-tappka.tex

\documentclass[11pt,a4paper]{article}

\usepackage[T1]{fontenc}
\usepackage[utf8]{inputenc}
\usepackage[czech]{babel}
\usepackage[sfdefault]{roboto}
\usepackage[scaled=0.96]{montserrat}
\renewcommand{\familydefault}{Roboto-TLF}
\newcommand{\headfont}{\fontfamily{Montserrat-TLF}\selectfont}

\usepackage[margin=2.2cm,headheight=14pt,footskip=1.1cm]{geometry}
\usepackage{tikz}
\usepackage{booktabs}
\usepackage{tabularx}
\usepackage{array}
\usepackage{enumitem}
\usepackage{xcolor}
\usepackage{tcolorbox}
\usepackage{fancyhdr}
\usepackage{titlesec}
\usepackage{microtype}
\usepackage[unicode,hidelinks]{hyperref}
\usepackage{xurl}  % dovolí zalomit dlouhé URL

% --- Barvy dle DESIGN.md -------------------------------------------------
\definecolor{tapred}{HTML}{B31B1B}      % primary
\definecolor{tapbrown}{HTML}{2C1A1D}    % foreground / secondary
\definecolor{tapwhite}{HTML}{FCFFF7}    % background
\definecolor{tapmuted}{HTML}{F2F4EE}    % muted
\definecolor{tapaccent}{HTML}{F5EFE3}   % accent
\definecolor{tapline}{HTML}{DFDCD4}     % border
\definecolor{tapgrey}{HTML}{6B6165}     % muted foreground
\definecolor{tapwarn}{HTML}{8A5200}     % warning-strong
\definecolor{tapwarnbg}{HTML}{FBF3E8}
\definecolor{tapok}{HTML}{1F6B34}       % success-strong

\color{tapbrown}
\pagecolor{tapwhite}

% --- Hlavička a patička --------------------------------------------------
\pagestyle{fancy}
\fancyhf{}
\fancyhead[L]{\footnotesize\color{tapgrey}Tappka \textbar{} Zálohování a obnova dat}
\fancyhead[R]{\footnotesize\color{tapgrey}verze 2.0 \textbar{} 1.\,9.\,2026}
\fancyfoot[C]{\footnotesize\color{tapgrey}\thepage}
\renewcommand{\headrulewidth}{0.4pt}
\renewcommand{\headrule}{\color{tapline}\hrule height 0.4pt}

% --- Nadpisy -------------------------------------------------------------
\titleformat{\section}
  {\headfont\bfseries\large\color{tapred}}{\thesection.}{0.55em}{}
\titlespacing*{\section}{0pt}{12pt}{5pt}
\titleformat{\subsection}
  {\headfont\bfseries\normalsize\color{tapbrown}}{\thesubsection}{0.5em}{}
\titlespacing*{\subsection}{0pt}{10pt}{2pt}

% --- Boxy ----------------------------------------------------------------
\newtcolorbox{pozor}{
  colback=tapwarnbg, colframe=tapwarn, boxrule=0pt, leftrule=3pt,
  arc=3pt, left=11pt, right=11pt, top=7pt, bottom=7pt}
\newtcolorbox{klid}{
  colback=tapmuted, colframe=tapok, boxrule=0pt, leftrule=3pt,
  arc=3pt, left=11pt, right=11pt, top=7pt, bottom=7pt}
\newcommand{\boxhead}[2]{{\headfont\bfseries\color{#1}#2}\par\smallskip}

% --- Značka Tappka -------------------------------------------------------
% Vektorová obdoba public/icon.svg: zaoblený čtverec s písmenem T.
% Kresleno přímo, aby značka byla ostrá v tisku a měla přesné barvy z DESIGN.md.
\newlength{\logosize}
\setlength{\logosize}{21mm}
\newcommand{\tappkalogo}{%
  \begin{tikzpicture}[x=\logosize, y=\logosize]
    \fill[tapred, rounded corners=0.25\logosize] (0,0) rectangle (1,1);
    \node[anchor=base, inner sep=0pt, text=tapwhite] at (0.5,0.297)
      {\headfont\bfseries\fontsize{0.547\logosize}{0pt}\selectfont T};
  \end{tikzpicture}}

% --- Štítek stavu --------------------------------------------------------
\newcommand{\stitek}[3]{%
  \tcbox[on line, colback=#1, colframe=#1, boxrule=0pt, arc=2pt,
         left=6pt, right=6pt, top=2.5pt, bottom=2.5pt, nobeforeafter]%
  {\footnotesize\headfont\bfseries\color{#2}#3}}

% --- Tabulky -------------------------------------------------------------
\newcolumntype{R}[1]{>{\raggedright\arraybackslash}p{#1}}
\newcolumntype{L}[1]{>{\raggedright\arraybackslash\headfont\bfseries\small}p{#1}}
\renewcommand{\arraystretch}{1.22}
\setlength{\heavyrulewidth}{0.9pt}
\setlength{\lightrulewidth}{0.35pt}
\setlength{\emergencystretch}{3em}
\setlength{\parindent}{0pt}
\setlength{\parskip}{5pt}
\setlength{\tabcolsep}{5pt}
\newcommand{\hd}[1]{\headfont\bfseries\small\color{tapbrown}#1}
\setlist{itemsep=2pt, topsep=3pt, parsep=0pt}

% LaTeX standardně zvýhodňuje zlom stránky před \section (\@secpenalty).
% V kombinaci s \raggedbottom to trhá stránku i tam, kde je ještě místo;
% zlomy tady řídíme ručně přes \clearpage.
\makeatletter
\def\@secpenalty{0}
\makeatother
\widowpenalty=10000
\clubpenalty=10000

% Přehledová tabulka parametrů, stejná struktura pro všechny tři způsoby.
% Musí být makro s argumentem — tabularx čte tělo jako argument, takže
% ho nelze rozdělit mezi \begin a \end vlastního prostředí.
\newcommand{\parametry}[1]{%
  \begin{tabularx}{\textwidth}{@{}L{3.5cm}X@{}}\toprule #1 \bottomrule\end{tabularx}}

% =========================================================================
\begin{document}

% =========================================================================
% STRANA 1 — HLAVNÍ
% =========================================================================
\thispagestyle{empty}

\vspace*{34mm}

\tappkalogo

\vspace{13mm}

{\headfont\bfseries\fontsize{27}{31}\selectfont\color{tapbrown}Zálohování a~obnova dat\par}

\vspace{7mm}

{\headfont\large\color{tapred}Tappka\par}
\vspace{1mm}
{\color{tapgrey}\normalsize Studentský portál Tiimiakatemia Praha\par}

\vspace{9mm}
\noindent{\color{tapred}\rule{34mm}{2.5pt}}

\vspace{7mm}
{\color{tapgrey}Dokument předkládá tři způsoby zálohování dat aplikace Tappka\\
a~shrnuje požadavky na součinnost IT oddělení PEF.\par}

\vfill

\begin{tabularx}{\textwidth}{@{}L{3.5cm}X@{}}
\toprule
Verze a~datum & 2.0 \textbar{}~1.\,9.\,2026 \\
Určeno & IT oddělení PEF, vedení školy \\
Zpracoval & Vývojový tým Tappka \\
Stav & Návrh k~projednání \\
\bottomrule
\end{tabularx}

\vspace{8mm}


% =========================================================================
% STRANA 2 — ZPŮSOB I
% =========================================================================
\clearpage
\section{Způsob I --- automatická záloha u~poskytovatele}

\stitek{tapmuted}{tapok}{V PROVOZU}\hspace{6pt}%
\stitek{tapmuted}{tapgrey}{BEZ NÁKLADŮ NAVÍC}\hspace{6pt}%
\stitek{tapmuted}{tapgrey}{BEZ SOUČINNOSTI ŠKOLY}

\parametry{%
Rozsah & Databáze PostgreSQL (texty esejí, hodnocení, rezervace, docházka) \\
Frekvence & Jednou denně, automaticky \\
Doba uchování & 7 dní \\
Umístění & Supabase, region Curych --- tedy tatáž infrastruktura jako produkční data \\
Obnovu provádí & Vývojový tým Tappka, samoobslužně \\
Doba obnovy & Desítky minut, podle objemu dat \\
Dostupnost při obnově & Aplikace je po dobu obnovy nedostupná \\
}

\subsection{Co tento způsob řeší}

\begin{itemize}
  \item Omylem smazaný nebo přepsaný záznam, pokud je zjištěn do sedmi dní.
  \item Chybu při aktualizaci aplikace, která poškodí strukturu databáze.
  \item Technické selhání na straně databázového serveru.
\end{itemize}

\begin{pozor}
\boxhead{tapwarn}{Omezení}
\begin{enumerate}[leftmargin=1.3em, itemsep=4pt]
  \item \textbf{Nahrané soubory nejsou zálohovány.} Fotografie a~dokumenty leží v~odděleném
        úložišti, na které se záloha nevztahuje; databáze o~nich obsahuje jen odkazy, nikoli
        obsah. Smazaný soubor je nenávratně ztracen. Omezení platí u~poskytovatele obecně.
  \item \textbf{Sedmidenní lhůta je krátká.} Nedostatek zjištěný po prázdninách, po zkouškovém
        období nebo po delší nepřítomnosti vyučující:ho již nelze napravit.
  \item \textbf{Zálohy leží u~téhož poskytovatele jako originál} --- neposkytují proto ochranu
        proti jeho výpadku, ukončení služby ani ztrátě přístupu k~účtu.
  \item \textbf{Poskytovatel v~tomto tarifu neposkytuje smluvní garanci dostupnosti.} V~případě
        výpadku neexistuje vymahatelná lhůta obnovení služby.
\end{enumerate}
\end{pozor}

Způsob I pokrývá běžné provozní chyby, jako jediné opatření je však nedostatečný: nechrání
nahrané soubory a~nezajišťuje nezávislost na dodavateli. Obojí řeší způsob II.

\subsection{Proč aplikace neběží na serveru školy}

Platformu lze provozovat i~ve vlastní režii, na infrastruktuře školy; tuto variantu jsme
zvažovali. Podle dokumentace poskytovatele\footnote{\url{https://supabase.com/docs/guides/self-hosting}, ověřeno k~1.\,9.\,2026.} však ve vlastní režii přechází na provozovatele
\emph{údržba serveru, bezpečnostní aktualizace, správa databáze, dostupnost, monitoring
a~zálohy}.

\begin{itemize}
  \item \textbf{Vlastní instalace neobsahuje žádné zálohy} --- celý způsob I by bylo nutné
        vytvořit a~udržovat vlastními silami. Vlastní provoz by tedy riziko ztráty dat spíše
        zvýšil než snížil.
  \item \textbf{Provozní odpovědnost by nesla škola}, včetně výpadků mimo pracovní dobu;
        vývojový tým Tappky nemá kapacitu držet nepřetržitou pohotovost.
  \item \textbf{Náklady jsou nízké} --- přibližně 600 Kč měsíčně, tedy zlomek času, který by
        správa vlastní instalace vyžadovala.
  \item \textbf{Rozhodnutí je vratné} --- data jsou ve standardním PostgreSQL a~způsob II drží
        jejich kopii u~školy; přechod na vlastní provoz zůstává otevřený.
\end{itemize}

% =========================================================================
% STRANA 3 — ZPŮSOB II
% =========================================================================
\clearpage
\enlargethispage{\baselineskip}
\section{Způsob II --- nezávislá kopie na serveru školy}
\label{sec:zpusob2}

\stitek{tapaccent}{tapred}{NÁVRH --- PŘEDMĚT JEDNÁNÍ}\hspace{6pt}%
\stitek{tapmuted}{tapgrey}{VYŽADUJE SOUČINNOST IT PEF}

\parametry{%
Rozsah & \textbf{Databáze i~nahrané soubory} --- úplná kopie dat \\
Frekvence & Jednou týdně, automaticky \\
Doba uchování & 8 týdenních, 12 měsíčních a~3 roční body obnovy \\
Umístění & \textbf{Server školy v~České republice} \\
Provoz & Kontejnerizovaná služba (Docker), spouštěná plánovačem \\
Zabezpečení & Kontrolní součty pro ověření neporušenosti archivů \\
Ověřování & Čtvrtletní zkušební obnova s~písemným protokolem \\
}

\subsection{Princip}

Služba na serveru školy si jednou týdně sama stáhne úplnou kopii dat --- databázi i~nahrané
soubory. Po prvotním nastavení pracuje samostatně a~upozorní pouze při selhání.

\subsection{Co tento způsob řeší}

\begin{itemize}
  \item \textbf{Ztrátu nahraných souborů} --- jediné opatření, které fotografie a~dokumenty
        chrání.
  \item \textbf{Nedostatek zjištěný po delší době} --- historie sahá až tři roky zpět.
  \item \textbf{Výpadek či ukončení služby poskytovatele} --- data zůstávají v~držení školy.
\end{itemize}

\begin{klid}
\boxhead{tapok}{Obnovení provozu z~této zálohy}
Záloha je \textbf{úplná a~plnohodnotná} --- obsahuje veškerá data v~otevřeném formátu a~lze ji
načíst do libovolné databáze PostgreSQL, na serveru školy i~u~jiného poskytovatele.

\textbf{Zprovoznění však vyžaduje technickou znalost}: nainstalovat databázový server, zálohu
naimportovat a~nastavit přístupová práva. Nejde o~úkon pro netechnickou osobu. Jde však
o~\textbf{nejrozšířenější otevřený databázový systém}, takže potřebná odbornost je běžně
dostupná a~nejsme závislí na konkrétním dodavateli.

Uvedení celé aplikace do provozu jinde je odhadováno na jednotky dní; převažující část této
doby není technická, nýbrž organizační.
\end{klid}

\subsection{Požadovaná kapacita}

\textbf{Současné využití je přibližně 1~\% dostupné kapacity.} Kalkulace níže proto nepopisuje
dnešní potřebu, nýbrž \emph{strop} --- stav, kdy by aplikace vyčerpala veškerou kapacitu,
kterou má u~poskytovatele k~dispozici.

\begin{tabularx}{\textwidth}{@{}Xrr@{}}
\toprule
\multicolumn{1}{@{}l}{\hd{Položka}} & \hd{Dnes} & \hd{Při plné kapacitě} \\
\midrule
Kopie nahraných souborů včetně verzí & jednotky GB & 160 GB \\
Body obnovy databáze (23 celkem) & jednotky GB & 115 GB \\
Pracovní prostor pro ověření obnovy & 10 GB & 150 GB \\
\midrule
Součet & \textasciitilde{}\,15 GB & 425 GB \\
Rezerva na růst dat a~prodloužení historie & --- & 425 GB \\
\midrule
\textbf{Navrhovaná alokace} & & \textbf{1 TB} \\
\bottomrule
\end{tabularx}

Alokace 1~TB je dimenzována na mnoho let dopředu. \textbf{Vyhovuje nám však i~start s~menším
objemem} (např. 100~GB s~možností navýšení) --- zálohování začne fungovat okamžitě.

% =========================================================================
% STRANA 4 — ZPŮSOB III
% =========================================================================
\clearpage
\section{Způsob III --- lidsky čitelný export dat}

\stitek{tapmuted}{tapwarn}{KONCEPT}\hspace{6pt}%
\stitek{tapmuted}{tapgrey}{BEZ SOUČINNOSTI ŠKOLY}\hspace{6pt}%
\stitek{tapmuted}{tapgrey}{VÝVOJOVÁ KAPACITA}

\parametry{%
Rozsah & Eseje, reflexe, profily dovedností, záznamy o~trénincích a~další studijní záznamy, včetně příloh \\
Frekvence & Na konci semestru a~kdykoli na vyžádání \\
Doba uchování & Neomezená; archiv je běžný soubor \\
Umístění & Dle volby školy --- školní disk, školní cloud, přenosné médium \\
Formát & Archiv ZIP se složkami; dokumenty ve formátu PDF, přehled v~CSV \\
Předpoklady obnovy & Žádné --- postačuje běžný počítač \\
Stav & Koncept; není součástí požadavků na školu \\
}

\subsection{Princip}

Způsoby I a~II chrání data v~technické podobě a~oba předpokládají, že bude k~dispozici osoba
schopná obnovit databázi. Způsob III řeší situaci, kdy tento předpoklad neplatí.

Naprostá většina záznamů v~Tappce je text nebo strukturovaný přehled; obojí lze uložit jako
samostatný dokument. Navrhujeme doplnit do aplikace funkci, která vytvoří archiv v~této
struktuře:

\begin{tcolorbox}[colback=tapmuted, colframe=tapline, boxrule=0.6pt, arc=3pt,
                  left=11pt, right=11pt, top=8pt, bottom=8pt]
\ttfamily\small
tappka-2026-podzim/\\
\hspace*{1.4em}Nováková Jana/\\
\hspace*{2.8em}Eseje/2026-10-12 Reflexe projektu.pdf\\
\hspace*{2.8em}Eseje/2026-11-03 Kniha -- Lean Startup.pdf\\
\hspace*{2.8em}Profil dovedností.pdf\\
\hspace*{2.8em}Tréninky a~aktivity.pdf\\
\hspace*{1.4em}Svoboda Petr/\\
\hspace*{2.8em}\dots\\
\hspace*{1.4em}PREHLED.csv
\end{tcolorbox}

Rozsah exportu není pevně dán --- nové typy záznamů se do archivu doplní stejným způsobem.

\subsection{Co tento způsob řeší}

\begin{itemize}
  \item \textbf{Zachování studijních záznamů v~čitelné podobě} nezávisle na existenci
        aplikace i~databáze.
  \item \textbf{Přístup k~materiálům během výpadku} --- vyučující může číst i~hodnotit offline.
  \item \textbf{Žádost studující:ho o~kopii vlastních údajů} podle nařízení GDPR.
  \item \textbf{Doložení stavu dokumentace} za uzavřený akademický rok.
\end{itemize}

\begin{klid}
\boxhead{tapok}{Obnovení provozu z~tohoto exportu}
Archiv lze rozbalit na kterémkoli počítači a~otevřít běžnou čtečkou dokumentů; přehled
v~tabulkovém procesoru. \textbf{Není zapotřebí databáze, instalace ani technická součinnost.}

Jde o~nejrychleji dostupné provizorium: materiály jsou k~dispozici řádově do jedné hodiny
a~lze je číst, tisknout, hodnotit i~zasílat elektronickou poštou.
\end{klid}

\subsection{Stav}

Způsob III je konceptem, nikoli hotovou funkcí. Nevyžaduje žádné rozhodnutí ani součinnost ze
strany školy a~je uveden proto, aby byl model ochrany dat popsán úplně. Jeho realizace závisí
na vývojové kapacitě týmu.

\end{document}
